%! Author = chtompki
%! Date = 1/14/26
\documentclass[12pt]{amsart}
% Default margins are too wide all the way around. I reset them here
\setlength{\hoffset}{-.75in}
\setlength{\textwidth}{6.5in}
\setlength{\voffset}{-.5in}
\setlength{\textheight}{9.0in}
\usepackage{amsmath}
\usepackage{leftindex}
\usepackage{amssymb}
\usepackage{amsthm}
\usepackage{enumitem}
\usepackage{setspace}

\newenvironment{solution}
{\begin{proof}[Solution]}
{\end{proof}}

\title{Combinatorics\\Exam 1}
\author{Rob Tompkins}
\date{\today}

\begin{document}
    \maketitle
    \date{\today}
    \begin{onehalfspace}
        % Problem 1
        \noindent\textbf{1.} Answer the following using simple counting reasoning
        \begin{enumerate}[label=(\alph*)]
            \item How many different 7-digit numbers (non-negative integers) are there?
                  (Leading zeros are not allowed).
            \item How many even 7-digit numbers are there?
            \item How many 7-digit numbers are there with exactly one 7?
            \item How many 7-digit numbers read the same when you read themm forward or backward, i.e. are palindromic,
                  e.g. 6548456 is one such integer.
            \item How many 7-digit numbers are there that are divisible by 4?
            \item How many times is the digit 7 written when listing the numbers between 1 and 1,000,000?
        \end{enumerate}
        \begin{solution}
        \end{solution}

        % Problem 2
        \noindent\testbf{2.} How many 7-letter words can be constructed by using the 26 letters of the alphabet if
        each word contains 3, or 4 vowels?
        It is understood that there is no restriction on the numger of times a letter can be used in a word.
        \begin{solution}
        \end{solution}

    \end{onehalfspace}
\end{document}
